\subsection{A typed fibre dictionary}\label{fam:section:B.1}

\begin{proposition}\label{fam:src:B.1.1}\label{fam:st:B.1.1}
Let $E_\tau=\mathbb C/(\mathbb Z+\tau\mathbb Z)$, $P=[\mu]$, and $L_\mu=\mathcal O_{E_\tau}(P-O)$. The holomorphic map
\[
 [\xi,\eta]\longmapsto[\xi,e^{2\pi i\eta}]
\]
identifies $\mathbb C^2/\Pi_c\Lambda$ with the quotient of $L_\mu^\times$ by the lift
$A(\xi,W)=(\xi+6\mu,e^{2\pi i\beta_c}W)$. Engel's generator, in Definition 2.3 of the attached paper, is its inverse, with translation $-6P$ and multiplier $e^{-2\pi i\beta_c}$. The section $P=[\mu]$ is a generator of the Mordell--Weil group of the stated pointed Weierstrass surface.
\end{proposition}
\Uses{the hypotheses, notation and definitions stated here; \cref{fam:src:A.3.1}}

\begin{proof}
First quotient by columns two, three and four. Exponentiation kills column four; columns three and two then give
\[
 (\xi,W)\sim(\xi+1,W)\sim(\xi+\tau,e^{2\pi i\mu}W).
\]
For the odd theta function with periods $1,\tau$, the quotient $\theta(\xi-\mu)/\theta(\xi)$ is a meromorphic section of this line bundle: its factors are $1,e^{2\pi i\mu}$ and its divisor is $P-O$. Sending its divisor section to this quotient is the asserted line-bundle isomorphism. Quotienting the remaining first column gives exactly $A$. On sufficiently small simply connected charts the logarithm of $W$ is a holomorphic inverse; changes of logarithm are column-four periods. Thus the displayed map is a biholomorphism of the two quotient complex spaces, with its inverse explicitly defined by local logarithms. Choosing $A^{-1}$ changes the generator of the same infinite cyclic group and gives precisely Engel's declared inverse-translation convention.

For clarity the elliptic surface in this dictionary is the displayed cubic, rather than only a curve with the same $j$-value. Let $f_{\rm sc}$ be the branch satisfying
\[
 f_{\rm sc}^4=\frac{12t^3(t-1)}{g_2(\tau)},\qquad
 f_{\rm sc}^6=-\frac{8t^4(t-1)^2}{g_3(\tau)},\qquad
 f_{\rm sc}/t\longrightarrow i\sqrt3/\pi.
\]
The identities $g_2^3/(g_2^3-27g_3^2)=t$ make these equations compatible. Their zero and pole orders give a holomorphic branch on the period cover; the indicated cusp sign fixes it. The map $\xi\mapsto(f_{\rm sc}^2\wp_\tau(\xi),\frac12f_{\rm sc}^3\wp'_\tau(\xi))$ is a pointed isomorphism to
$y^2=x^3-3t^3(t-1)x+2t^4(t-1)^2$. Its factors are $f_{\rm sc}(g_1z)=-f_{\rm sc}/\tau$, $f_{\rm sc}(g_2z)=f_{\rm sc}/\tau$; therefore this isomorphism descends over $B^\circ$.

We first verify that $P=[\mu]$ gives a rational section of the resolved elliptic surface. On $B^\circ$ its cubic coordinates descend by the displayed scale and affine laws. At the two finite centres, $\mu_j$ is not a lattice point, while $f_{\rm sc}$ has orders two and one in the root parameter. Its cubic coordinates are therefore holomorphic and invariant, hence holomorphic in the base coordinate. At the cusp $\mu$ is holomorphic and bounded, and the convergent multiplicative expansion of $\wp$ gives meromorphic cubic coordinates with finite pole orders. Even when its exponential meets one, the pole is controlled by its finite vanishing order; it cannot be identically one because the two finite values of $\mu$ differ. The coordinates are thus meromorphic on $\mathbb P^1$ and rational. The closure of its graph on the proper resolved surface is an irreducible curve; its normalization maps finitely with degree one to the smooth base, so is the base itself. This gives the unique section of that surface extending $P$.

We also verify that the specified section is a Mordell--Weil generator using its full logarithm. A rational section of this cubic pulls back to a holomorphic section of the good elliptic family at both finite covering points: multiply its projective coordinates by a common power to remove poles, and use the nonsingular limiting Weierstrass equation. At the cusp put $v=1/t$, $X=v^2x,Y=v^3y$. Its equation is $Y^2=X^3-3(1-v)X+2(1-v)^2$. A section cannot pass through the node $(1,0)$: substitution of $X=1+O(v)$, $Y=O(v)$ leaves the coefficient $+1$ of $v$ on the left-hand side of the defining equation. At every other central point the multiplicative uniformization is a local isomorphism. Explicitly, at $v=0$ it is
\[
 X=1+\frac{12U}{(U-1)^2},\qquad
 Y=\frac{12\sqrt3 U(U+1)}{(U-1)^3},\qquad U=e^{2\pi i\xi},
\]
with inverse $U=(Y/(X-1)+\sqrt3)/(Y/(X-1)-\sqrt3)$ away from $O$; a local parameter at $O$ is $X/Y$. The convergent lattice expansion of $\wp$ gives the same holomorphic map for small $v$. Hence a section has a holomorphic unit $U(v)$ and a bounded logarithm at the cusp. Lifting the section on the simply connected cover gives a global logarithm $\nu$, with this cusp branch. Its only remaining ambiguity is an integer, since a nonzero multiple of $\tau$ is unbounded.

Its two affine laws have integral defects:
\[
 \nu(g_1z)=-\nu/\tau+\alpha_1+\beta_1\tau(g_1z),\quad
 \nu(g_2z)=\nu/\tau+\alpha_2+\beta_2\tau(g_2z).
\]
Composing them and using cusp invariance gives $\alpha_2=-\beta_1$, $\beta_2=\alpha_1+\beta_1$. Put $n=-\alpha_1-2\beta_1$ and $a=\alpha_1+\beta_1$. Substitution shows that $n\mu+a$ has exactly these defects. Their difference is a global section of the first-coordinate sheaf $\mathcal O_B(-p_0)$ from \ref{fam:src:A.3.1}, so is zero. Every rational section is therefore $nP$. The same cusp argument and the distinct fixed values of $\mu$ show that $P$ has infinite order. This proves $\operatorname{MW}(S)=\mathbb ZP$. \end{proof}

\begin{proposition}\label{fam:src:B.1.2}\label{fam:st:B.1.2}
The exact set of nonzero parameters for which every ordinary fibre and both normalized finite root fibres of Engel's period-reference quotient are compact complex two-tori is
\[
 0<|u|<e^{-2\pi M}.
\]
This is the parameter domain for the compact family with both smooth finite root models. The condition on ordinary fibres alone at equality is determined by whether a maximizing point of $D_{\rm part}$ is ordinary; this boundary issue is treated in Part C. For an independent reference, the radius is $|\lambda_0|e^{-2\pi M}$.
\end{proposition}
\Uses{the hypotheses, notation and definitions stated here; the explicit constructions and calculations in the proof}

\begin{proof}
The function $N(\xi,W)=|W|\exp(2\pi\operatorname{Im}\mu\,\operatorname{Im}\xi/\operatorname{Im}\tau)$ is invariant under both line-bundle period identifications. The inverse lift multiplies it by $e^{2\pi D_c}$. More directly the real determinant of the four period columns is a nonzero sign times $(\operatorname{Im}\tau)D_c$. Thus the four columns generate a full lattice exactly when $D_c\ne0$, and then the quotient is a compact torus by a closed fundamental parallelepiped. At $D_c=0$ their real span has dimension three; a rank-four lattice in $\mathbb R^4$ cannot result, so the quotient cannot be a compact complex two-torus. The function $D_{\rm part}$ descends continuously to the connected base with the cusp removed and tends to $-\infty$ at the cusp. Its image is therefore $(-\infty,M]$: the maximum is attained, and the intermediate value theorem fills all smaller values. Consequently $D_c$ is nonzero at every base point exactly when $\operatorname{Im}c<-M$. Since $|u|=e^{2\pi\operatorname{Im}c}$, the assertion follows, with the indicated reference scaling.
\end{proof}

\begin{proposition}\label{fam:src:B.1.3}\label{fam:st:B.1.3}
For the original affine logarithm $\mu$, we have $\wp_\tau(\mu)=0$ identically. Equivalently the generator $P=[\mu]$ of \ref{fam:src:B.1.1} is one of the two sections $P_*=(0,\sqrt2t^2(t-1))$ and $-P_*$. This assertion does not use a numerical evaluation of $\mu$ or of $\beta$.
\end{proposition}
\Uses{the hypotheses, notation and definitions stated here; \cref{fam:src:B.1.1}}

\begin{proof}
We give the intersection calculation that identifies the generator. At $t=1$, put $u=t-1$ and blow up the origin. In the $u$-chart $x=uX,y=uY$, the strict equation is
$Y^2=uX^3-3(1+u)^3X+2(1+u)^4$. Its exceptional curve is smooth, and $P_*$ meets it at $(X,Y)=(0,\sqrt2)$. In the other chart $u=xU,y=xY$, the equation is $Y^2=x-3(1+xU)^3U+2(1+xU)^4U^2$. The implicit function theorem solves for $x$ at the common point; the two fibre branches $U=0,x=0$ meet with intersection multiplicity two. These charts cover the tangent conic and resolve the singularity. The two components have multiplicity one; the one met by $P_*$ is the nonidentity component $D$, with $D^2=-2$ from intersection with the full fibre.

At $t=0$, two blow-ups in the successive $t$-charts give
\[
 Y^2=tX^3+3t^2(1-t)X+2t^2(1-t)^2,
\quad
 Y_2^2=t^2X_2^3+3t(1-t)X_2+2(1-t)^2.
\]
The second equation has two smooth multiplicity-one components $E_{2,\pm}$ at $Y_2=\pm\sqrt2$, and $P_*$ meets one of them. The remaining singular point, in the second $X$-chart, has
\[
 V^2=H(X,T),\quad H=X^2T+(2+3X)T^2-(4X+3X^2)T^3+2X^2T^4,
 \qquad t=XT.
\]
Let $T_0(X)$ solve $H_T(X,T_0)=0$. It has $T_0=-X^2/4+O(X^3)$. Taylor's formula gives
$H-H(X,T_0)=(T-T_0)^2h_0$, $h_0(0,0)=2$, and $H(X,T_0)=X^4e(X)$, $e(0)=-1/8$. These identities follow by substitution in the displayed polynomial, not by discarding its higher terms. With $Z=(T-T_0)\sqrt{h_0}$, $a=V-Z$, $b=(V+Z)/e(X)$, the equation becomes exactly $ab=X^4$ through a local analytic isomorphism. Its resolution has charts
\[
 a=U_i^{i+1}V_i^i,\quad b=U_i^{3-i}V_i^{4-i},\quad X=U_iV_i
 \quad(0\le i\le3),
\]
with changes $U_{i+1}=1/V_i,V_{i+1}=U_iV_i^2$. They give a chain of three $(-2)$-curves. These are the charts of the two successive point blow-ups resolving $ab=X^4$, so the resolution is proper. Retaining $t=XT$ in the comparison gives multiplicities $2,3,2$ on that chain. The two new multiplicity-one components meet its ends; the first exceptional curve meets its middle component, and the original strict fibre meets the other end of that first exceptional curve. All these statements follow by substituting the displayed monomials and their orders; the only fibre components of multiplicity one are the strict component and $E_{2,\pm}$.

Order the six nonidentity components by the graph
$\Theta_1-\Theta_2-\Theta_3-\Theta_4$ and $\Theta_2-\Theta_5-\Theta_6$, with $E_{2,\pm}=\Theta_4,\Theta_6$. Their intersection matrix is $-C$, where
\[
 C=\begin{pmatrix}2&-1&0&0&0&0\\-1&2&-1&0&-1&0\\0&-1&2&-1&0&0\\0&0&-1&2&0&0\\0&-1&0&0&2&-1\\0&0&0&0&-1&2\end{pmatrix},\quad
 C^{-1}=\frac13\begin{pmatrix}6&9&6&3&6&3\\9&18&12&6&12&6\\6&12&10&5&8&4\\3&6&5&4&4&2\\6&12&8&4&10&5\\3&6&4&2&5&4\end{pmatrix}.
\]
Multiplication verifies this inverse. Self-intersections follow as well from the complete multiplicity vector and $F\cdot\Theta_i=0$.

The form $dx\wedge dt/y$ is holomorphic and nonvanishing on the finite resolved charts. In the last charts its expression is a unit times $dX\wedge da/a=-dU_i\wedge dV_i$; the earlier blow-up equations cancel the apparent denominators at $y=0$. At infinity it is $-v^{-1}dX\wedge dv/Y$ and has one fibre pole. Thus $K_S=-F$. Adjunction gives $Q^2=-1$ for every section $Q$, since $F\cdot Q=1$. A section meets a unique multiplicity-one component transversely at a smooth fibre point: the pullback of the full base equation to the section has order one, excluding all other possibilities.

Let $v_Q$ be the intersection vector with $\Theta_i$ and let $d_Q=Q\cdot D$. The possibilities are $v_Q=0,e_4,e_6$ and $d_Q=0,1$. Define the rational divisor class
\[
 \varphi(Q)=Q-O-(Q\cdot O+1)F+\sum\Theta_i(C^{-1}v_Q)_i+\tfrac12d_QD.
\]
It is orthogonal to $O,F,\Theta_i,D$ by direct multiplication. The chord-and-tangent rational function on the generic cubic has divisor $(Q)+(R)-(Q+R)-(O)$; extending it to the surface adds only vertical components. Applying the displayed orthogonal projection therefore proves $\varphi(Q+R)=\varphi(Q)+\varphi(R)$. Expansion gives
\[
 h(Q)=-\varphi(Q)^2=2+2Q\cdot O-v_Q^tC^{-1}v_Q-\tfrac12d_Q.
\]
For $Q\ne O$, the nonnegative intersection $Q\cdot O$ gives $h(Q)\ge2-4/3-1/2=1/6$. The visible $P_*$ is disjoint from $O$ everywhere, and the explicit blow-up charts place it on both nonidentity components. Hence $h(P_*)=1/6$. The logarithm argument in \ref{fam:src:B.1.1} already proved $P_*=nP$ for an integer $n\ne0$, without identifying the visible section. Additivity gives $1/6=n^2h(P)\ge n^2/6$, so $n=\pm1$. The pointed cubic map in \ref{fam:src:B.1.1} has $x(P)=f_{\rm sc}^2\wp_\tau(\mu)$ with $f_{\rm sc}\ne0$ on the ordinary locus. This proves the asserted zero identity.
\end{proof}

\begin{proposition}\label{fam:src:B.1.4}\label{fam:st:B.1.4}
The line-bundle comparison of the period family with Engel's family over $B^\circ$ transports its inverse period lift to a global section of $\operatorname{Hom}(t_{6P}^*M_S,M_S)$ with divisor $F_\infty$. Take this transported section for $c=0$ as the reference linearization. Then $\lambda_0=1$ exactly and $u=e^{-2\pi ic}=w^2$. For another independently fixed reference the precise constant is its scalar ratio with this section. With the cusp branch fixed in Part C and $f_{\rm sc}/t\to i\sqrt3/\pi$, the section is $P=(0,-\sqrt2t^2(t-1))$.
\end{proposition}
\Uses{the hypotheses, notation and definitions stated here; \cref{fam:src:B.1.3}; \cref{fam:src:B.1.1}; \cite[Lemma 2.1]{Engel2026}; \cite[Corollary 2.2]{Engel2026}}

\begin{proof}
The zero identity $\wp_\tau(\mu)=0$, proved without numerical input in \ref{fam:src:B.1.3}, identifies $P$ with the signed visible section. With the cusp branch in Part C and the specified scale, substitution of $U=-(5+2\sqrt6)$ in the nodal uniformization of \ref{fam:src:B.1.1} gives $Y=-\sqrt2$, so $P=(0,-\sqrt2t^2(t-1))$.

To compare the line bundles as families, use their common divisor $P-O$ and rigidify at the zero section. Put $F_\theta(\xi)=\theta(\xi-\mu)/\theta(\xi)$, $r_\theta=\theta(-\mu)/\theta'(0)$, and let $s_{P-O}$ be the canonical meromorphic divisor section. The map
\[
 [\xi,W]\longmapsto\frac{2r_\theta}{f_{\rm sc}}
        W\frac{s_{P-O}(\xi)}{F_\theta(\xi)}
\]
is a holomorphic line-bundle isomorphism, including at $P,O$ after cancellation of the equal divisors. At $O$ it sends the chosen frame to $dx/y$, since $F_\theta\sim r_\theta/\xi$ and $x/y\sim-\xi/f_{\rm sc}$ with the corresponding canonical divisor frame. Two such maps that agree at $O$ differ by a holomorphic scalar on each compact elliptic fibre and hence agree everywhere. This proves descent under the actual period identifications.

Engel's \cite[Lemma 2.1]{Engel2026}, which we have read together with its proof, identifies $\operatorname{Hom}(t_{6P}^*M_S,M_S)$ with $\pi_S^*\mathcal O_B(1)$. The lemma applies because the preceding pointed cubic map and Mordell--Weil calculation identify his $(S,O,P)$ with ours, including the chosen sign. To determine which section is transported, it suffices to compute its order at one generic smooth point of each special fibre, since a section of this pulled-back line is a base section. At a finite centre write $6\mu=p+r\tau+\delta$ with $(p,r)=(4,-2),(3,-3)$ and $\delta/f_{\rm sc}$ holomorphic; this divisibility follows by linearizing the affine laws, whose homogeneous orders are two and one. Evaluate at $\xi=f_{\rm sc}z$ for generic fixed $z$. The factor $(2r_\theta/f_{\rm sc})/F_\theta(f_{\rm sc}z)$ tends to $2z$, a unit. Restoring the lattice shift $p+r\tau$ at the target changes the fibre multiplier to $e^{-2\pi i(\beta_c-r\mu)}$, whose finite limits are $e^{-2\pi i\kappa_1}$ and $e^{-2\pi i\kappa_2}$. The source and target theta ratios remain units. Therefore the base section has order zero at both finite points. At the cusp $f_{\rm sc}$ has a simple pole and $r_\theta$ has a nonzero limit; the two factors $2r_\theta/f_{\rm sc}$ cancel between source and target. The remaining multiplier is
$e^{-2\pi i\beta_c}=q e^{-2\pi i(\beta_c+\tau)}$, of exact order one. Thus the transported section has divisor $F_\infty$. Engel's \cite[Corollary 2.2]{Engel2026} now identifies it with a nonzero scalar multiple of his reference, since two base sections with the same divisor have constant ratio. Changing $\beta$ by $c$ multiplies it exactly by $e^{-2\pi ic}$. With our stated period-reference normalization, this scalar is one. For any independently chosen reference $\phi_{\rm ref}$, the scalar is instead the precisely defined ratio $\lambda_0=\phi_{\rm part}/\phi_{\rm ref}$, evaluated on any ordinary fibre. Scaling the reference by $a\in\mathbb C^*$ changes $\lambda_0$ to $\lambda_0/a$; no scalar value is invariant under that declared choice.
\end{proof}


\section{The normalized periods, quantitative bounds, and the boundary}\label{fam:section:C}

\subsection{A theta expression for the missing period}\label{fam:section:C.1}
We use the odd Jacobi function
\[
 \vartheta_1(v,\tau)=2\sum_{n\geq0}(-1)^n
 e^{\pi i\tau(n+1/2)^2}\sin((2n+1)\pi v)
\]
and the Dedekind function
\(\eta_D(\tau)=e^{\pi i\tau/12}\prod_{n\geq1}(1-e^{2\pi in\tau})\).
The subscript distinguishes this function from the invariant integral class.
We choose the cusp lift reached from \(\tau=iT\), \(T>1\), on which
\(\mu=1/2-iy(T)\), \(0<y(T)<T/2\). Put
\(a=5+2\sqrt6\) and \(\rho=e^{\pi i/3}\).

\begin{proposition}\label{fam:src:C.1.1}\label{fam:st:C.1.1}
For the period section identified in Statement \ref{fam:src:B.1.3}, the particular solution
with \((\beta^{\rm part}+\tau)(p_0)=0\) is
\begin{equation}\label{fam:eq:cd-beta}
 \beta^{\rm part}=\frac{6i}{\pi}\operatorname{Log}
 \frac{\vartheta_1(\mu,\tau)}{2\sqrt3\,\eta_D(\tau)}.
\end{equation}
Here the logarithm is continued on the simply connected uniformizing
half-plane from its cusp expansion. Its special values are
\[
 \mu_0=\frac12-\frac{i\log a}{2\pi},\qquad
 \mu(z_1)=\frac{2-\rho}{3},\qquad \mu(z_2)=\frac{1-i}{2},
\]
\[
 \beta^{\rm part}(z_2)=i\left(\frac32-\frac3\pi\log6\right),
 \qquad \operatorname{Re}\beta^{\rm part}(z_1)=-\frac12.
\]
In particular \(e^{2\pi i\mu_0}=-a\).
\end{proposition}
\Uses{the hypotheses, notation and definitions stated here; \cref{fam:src:B.1.3}}

\begin{proof}
Statement \ref{fam:src:B.1.3} identifies the elliptic point represented by \(\mu\) with
one of the sections \((0,\pm\sqrt2t^2(t-1))\). The uniformization map
\((v,\tau)\mapsto(f^2\wp_\tau(v),f^3\wp'_\tau(v)/2)\) therefore gives
\(\wp_\tau(\mu)=0\). No zero of \(\wp\) is a lattice point. Thus the numerator
of \eqref{fam:eq:cd-beta} is nonzero; the denominator is nonzero as well.
A holomorphic logarithm exists on the simply connected cover.

For completeness the cusp value follows directly from the normally
convergent Fourier expansion, rather than from a quoted Eichler--Zagier
constant:
\begin{equation}\label{fam:eq:cd-wp}
 \frac{\wp_\tau(v)}{\pi^2}
 =\csc^2(\pi v)-\frac13+
 8\sum_{n\geq1}\frac{nq^n}{1-q^n}(1-\cos(2\pi nv)),
 \qquad q=e^{2\pi i\tau}.
\end{equation}
This identity follows by differentiating the logarithm of the Jacobi
product twice and using
\(E_2=1-24\sum nq^n/(1-q^n)\). Boundedness of the chosen elliptic logarithm
at the cusp permits passage to the limit. We obtain
\(\sin^2(\pi\mu_0)=3\), equivalently
\(u^2+10u+1=0\), \(u=e^{2\pi i\mu_0}\).
On the chosen imaginary-axis branch \(\operatorname{Im}\mu<0\), so \(u=-a\).
One can see the branch without a numerical convention: on the rectangular
elliptic curve with \(\tau=iT\), \(T>1\), \(\wp\) is real on the segment
\(1/2-iy\), \(0\leq y\leq T/2\). Its endpoint values are the largest and
middle real roots of its cubic, respectively positive and negative.
Its derivative has no zero in the open segment, since the zeros of
\(\wp'\) are the three nonzero half-periods. There is exactly one zero
there; this gives the indicated continuation. Consequently
\(2\sin(\pi\mu_0)=2\sqrt3\).

Put \(R=\vartheta_1/\eta_D\). Poisson summation for the defining Gaussian
series and its index shift give
\[
 \begin{aligned}
 R(v,\tau+1)&=e^{\pi i/6}R(v,\tau),\\
 R(v/\tau,-1/\tau)&=-i e^{\pi iv^2/\tau}R(v,\tau),\\
 R(1-v,\tau)&=R(v,\tau),\qquad R(v+1,\tau)=-R(v,\tau).
 \end{aligned}
\]
Inserting the actual transformed arguments gives
\[
 \begin{aligned}
 R(\mu\circ g_1,\tau\circ g_1)
 &=e^{-\pi i/3}e^{\pi i(1-\mu)^2/\tau}R(\mu,\tau),\\
 R(\mu\circ g_2,\tau\circ g_2)
 &=i e^{\pi i\mu^2/\tau}R(\mu,\tau).
 \end{aligned}
\]
Multiplication of the logarithms by \(6i/\pi\) gives exactly
\(2-6(1-\mu)^2/\tau\) and \(-3-6\mu^2/\tau\), modulo constants in
\(12\mathbb Z\). These constants vanish: each defect is constant on the
connected cover, and evaluation at the corresponding fixed point makes
both its two sides zero. The fixed-point equations themselves give
\(\mu(z_1)=(2-\rho)/3\) and \(\mu(z_2)=(1-i)/2\).
The Jacobi product has cusp expansion
\(R(\mu,\tau)=2q^{1/12}\sin(\pi\mu)(1+O(q))\).
Equation \eqref{fam:eq:cd-beta} consequently gives
\(\beta^{\rm part}+\tau\to0\). The difference from the specified particular
solution is an invariant holomorphic function extending across all three
points of \(B\), hence a constant, and the cusp normalization makes it zero.

At \(\tau=i\), the half-period translation formula gives
\(\vartheta_1((1-i)/2,i)=e^{\pi/4}\vartheta_3(0,i)\).
The identities \(2\eta_D^3=\vartheta_2\vartheta_3\vartheta_4\),
\(\vartheta_2(0,i)=\vartheta_4(0,i)=2^{-1/4}\vartheta_3(0,i)\)
give \(R=e^{\pi/4}\sqrt2\), proving the displayed value of \(\beta\).
At \(\rho\), \(q\) is negative real, \(\mu=1/2-i/(2\sqrt3)\), and the
product for \(R\) has phase \(\pi/12\), with all its remaining factors
positive real. This gives \(\operatorname{Re}\beta=-1/2\).
\end{proof}

\begin{proposition}\label{fam:src:C.1.2}\label{fam:st:C.1.2}
The following decimal evaluations of the exact expressions in \ref{fam:src:C.1.1} agree
in two independent evaluations through the displayed twenty decimal places:
\[
\begin{array}{c|l}
\mu_0&0.5-0.36485183191106784792\,i\\
\mu(z_1)&0.5-0.28867513459481288225\,i\\
\mu(z_2)&0.5-0.5\,i\\
\beta^{\rm part}(z_1)&-0.5-1.27101206057537346287\,i\\
\beta^{\rm part}(z_2)&-0.21100425815613411418\,i
\end{array}
\]
These are evaluations of explicit convergent expressions, with working
precision ninety decimal digits. The formula \eqref{fam:eq:cd-beta}, together
with \(j(\tau)=1728t\) and \(\wp_\tau(\mu)=0\), determines the periods
throughout a fundamental domain.
\end{proposition}
\Uses{the hypotheses, notation and definitions stated here; \cref{fam:src:C.1.1}; \cref{fam:eq:cd-beta}}

\begin{proof}
The executable \texttt{cd\_numerics.py} evaluates \(\mu_0\) first by its
algebraic logarithm and then by solving
\(\pi^2\csc^2(\pi v)-\pi^2/3=0\). For \(\beta\), its first calculation uses
the Jacobi theta series and the eta product; its second uses
\[
 R(v,\tau)=2e^{\pi i\tau/6}\sin(\pi v)
 \prod_{n\geq1}(1-2q^n\cos(2\pi v)+q^{2n}).
\]
The product is truncated after 120 factors. At the two specified points
\(|q|\leq e^{-\pi\sqrt3}\) and \(|\operatorname{Im}v|\leq1/2\), so the
omitted logarithm is bounded by
\(\sum_{n>120}(2|q|^n\cosh\pi+|q|^{2n})/(1-2|q|^n\cosh\pi-|q|^{2n})\),
which is less than \(10^{-280}\). The measured discrepancies of the
ninety-digit calculations are below \(10^{-88}\). We do not promote floating-point rounding to an interval certificate.
A further sixty-digit interval computation, \texttt{cd\_special\_intervals.py},
evaluates the real positive product at $\rho$, including its forty-factor
tail, and the exact logarithmic expression at $i$. It encloses both
imaginary parts within intervals narrower than $3\cdot10^{-59}$ and
therefore certifies the displayed twenty decimal places. It also proves
$-0.2<\operatorname{Im}(\kappa_1^{\rm part}-\kappa_2^{\rm part})<-0.1$.
The exact formulas, product-tail bounds, independent algorithms, and
full outputs accompany the decimals. The special values of \(\mu\) are also checked directly in
the fixed-point equations, an exact algebraic calculation.
\end{proof}

\subsection{Bounds for the sharp parameter disc}\label{fam:section:C.2}
\begin{proposition}\label{fam:src:C.2.1}\label{fam:st:C.2.1}
For the normalization of \ref{fam:src:C.1.1}, the exact admissibility maximum is finite
and satisfies
\[
 -\frac3\pi\log6\ \leq M<2.
\]
More precisely the proof below gives \(M<1.635\). Consequently the sharp
radius \(R_w=e^{-\pi M}\) satisfies
\(e^{-2\pi}<R_w\leq216\), and for the reference linearization with
\(u=w^2\), \(e^{-4\pi}<R_u\leq46656\).
\end{proposition}
\Uses{the hypotheses, notation and definitions stated here; \cref{fam:src:C.1.1}; \cref{fam:eq:cd-beta}; \cref{fam:eq:cd-norm}; \cref{fam:eq:cd-wp}; \cref{fam:src:C.2.3}}

\begin{proof}
Equation \eqref{fam:eq:cd-beta} gives
\begin{equation}\label{fam:eq:cd-norm}
 D_{\rm part}=\frac6\pi\log\frac{\mathcal N(\mu,\tau)}{2\sqrt3},
 \qquad
 \mathcal N(v,\tau)=e^{-\pi(\operatorname{Im}v)^2/\operatorname{Im}\tau}
 |R(v,\tau)|.
\end{equation}
The theta transformation formulas prove that \(\mathcal N\) is unchanged
by integral period translations of \(v\), sign change of \(v\), and the
simultaneous modular change of \(v,\tau\). We may therefore use the
comparison isomorphism \(\mathbb C/(\mathbb Z+\tau\mathbb Z)\to
\mathbb C/(\mathbb Z+\tau'\mathbb Z)\), \([v]\mapsto[v/(c\tau+d)]\), to
place \(\tau\) in the standard modular domain and then choose
\(|\operatorname{Im}v|\leq T/2\), \(T=\operatorname{Im}\tau\geq\sqrt3/2\).
Both operations preserve the equation \(\wp(v)=0\) and \eqref{fam:eq:cd-norm}.

For \(\sqrt3/2\leq T\leq2\), the theta series and
\(\prod(1-r^n)\geq\exp(-r/(1-r)^2)\) imply
\[
 \mathcal N\leq
 2e^{\pi/6}\left(1+\frac{e^{-\pi\sqrt3/2}}
 {1-e^{-3\pi\sqrt3/2}}\right)
 \exp\frac{e^{-\pi\sqrt3}}{(1-e^{-\pi\sqrt3})^2}
 <3.615.
\]
Indeed the absolute value of the \(n\)-th sine term, including the
Gaussian factor, is at most
\(2\exp(-\pi T(n+1/2-|\operatorname{Im}v|/T)^2)\), which is at most
\(2e^{-\pi Tn^2}\); and \(n^2\geq1+3(n-1)\) for \(n\geq1\).

For \(T\geq2\), a zero of \(\wp\) in this period strip has
\(|\operatorname{Im}v|<1\). To prove this, assume its absolute value is at
least one and apply \eqref{fam:eq:cd-wp}. With \(r=e^{-4\pi}\), the absolute
value of the terms other than \(-1/3\) is at most
\[
 \frac1{\sinh^2\pi}+8\left\{
 \frac r{(1-r)^3}+\frac1{2(1-r)}
 \left(\frac{e^{-2\pi}}{(1-e^{-2\pi})^2}
 +\frac{e^{-6\pi}}{(1-e^{-6\pi})^2}\right)\right\}<0.016<\frac13.
\]
Here \(|\sin\pi v|\geq\sinh\pi\), and
\(|q|^n\cosh(2\pi n\operatorname{Im}v)\leq
(e^{-\pi Tn}+e^{-3\pi Tn})/2\). This contradicts \(\wp(v)=0\).
The product expression for \(R\) now gives the uniform bound
\[
 \mathcal N\leq2\cosh\pi\,e^{-\pi/3}
 \exp\left(\frac{2\cosh(2\pi)r}{1-r}+\frac{r^2}{1-r^2}\right)
 <8.151.
\]
Insertion in \eqref{fam:eq:cd-norm} proves the stated upper bound. The script
\texttt{cd\_bounds.py} independently evaluates these monotone bounds by
fifty-digit interval arithmetic: the larger resulting bound for \(D\) is
contained, after outward rounding for display, in
\[
 [1.63421874605891012363,\;1.63421874605891012365].
\]
The output file retains the full fifty-digit endpoints.
This encloses the coarse upper bound; the exact maximum is proved in \ref{fam:src:C.2.3}.

At the cusp \(D_{\rm part}=-T+o(1)\), so it tends to minus infinity;
continuity on the rest of the compact orbifold base gives attainment of
its maximum. At \(z_2\), \ref{fam:src:C.1.1} gives
\(D_{\rm part}=-(3/\pi)\log6\), proving the lower bound. Addition of \(c\)
adds \(\operatorname{Im}c\) to \(D\); a zero is forced by continuity
between a maximum and the cusp whenever \(\operatorname{Im}c\geq-M\).
For \(\operatorname{Im}c<-M\), the real determinant \(TD\) is everywhere
nonzero, proving the exact asserted form of the admissible disc.
\end{proof}

\begin{proposition}\label{fam:src:C.2.2}\label{fam:st:C.2.2}
The point \(t=1\) is a strict local maximum of \(D_{\rm part}\).
A double-precision calculation on 5151 points of the half modular domain
\(0\leq\operatorname{Re}\tau\leq1/2\), \(|\tau|\geq1\),
\(\operatorname{Im}\tau\leq3\), has largest sampled value at \(\tau=i\).
This finite computation does not establish the global maximum.
\end{proposition}
\Uses{the hypotheses, notation and definitions stated here; the explicit constructions and calculations in the proof}

\begin{proof}
We give a quantitative neighbourhood, which will also specify the local
region needed for a global interval cover. Put
\[
 h=v-(1-\tau)/2,\quad x=\tau-i,\quad
 G(h,\tau)=\wp_\tau((1-\tau)/2+h),\quad
 H(h,\tau)=\log\frac{\vartheta_3(h,\tau)}{\eta_D(\tau)}.
\]
The half-period theta identity gives exactly
$\log\mathcal N=\operatorname{Re}H-\pi(\operatorname{Im}h)^2/T$.
On the complex polydisc $|h|\leq1/5$, $|x|\leq1/10$, the theta series
has $|\vartheta_3-1|<1$, so the logarithm is holomorphic. There are no
poles of $G$ there: for every integer lattice point its imaginary
difference from $(1-\tau)/2+h$ has absolute value at least
$0.9/2-0.2>0$. Both functions are even in $h$.

The following uniform bounds use only positive Gaussian sums:
\[
 |G|<20,\qquad |H|<0.6.
\]
For explicit verification put
$A_k=\sum_{n\geq1}n^k e^{-0.9\pi n^2}\cosh(0.4\pi n)$,
$a_0=2A_0$, $a_1=4\pi A_1$, $a_2=8\pi^2A_2$, and $r=e^{-1.8\pi}$.
The bounds used are
\[
 \begin{aligned}
 |G|&\leq\frac{a_2}{1-a_0}+\frac{a_1^2}{(1-a_0)^2}
 +\frac{\pi^2}3\left(1+\frac{24r}{(1-r)^3}\right)<18.373,\\
 |H|&\leq-\log(1-a_0)+\frac{1.1\pi}{12}+
 \frac r{(1-r)^2}<0.547.
 \end{aligned}
\]
The derivative identity behind the first bound is
$G=-\partial_h^2\log\vartheta_3-\pi^2E_2/3$.
The script \texttt{cd\_local\_bound.py} evaluates these bounds by sixty-digit
interval arithmetic. It truncates the positive sums at $n=10$ and bounds
the remaining terms by geometric sums using
$n^2\geq121+23(n-11)$; the polynomial factors $n^k$ are summed explicitly.

At $(0,i)$ the theta heat equation gives
\[
 G=0,\quad G_{h^2}=-g_2(i)/4,\quad G_x=id,
 \qquad H=\log\sqrt2,\quad H_x=0,\quad H_{h^2}=-\pi/2.
\]
Subscripts $h^2$ denote power-series coefficients. The first two follow
also from $\wp''=6\wp^2-g_2/2$. Here $d$ is the real positive derivative coefficient of
$e_2(\tau)=(\pi^2/3)(\vartheta_2(0,\tau)^4-
\vartheta_4(0,\tau)^4)$ at $i$. Differentiating the convergent theta
series, with explicit Gaussian tails, gives $10.03<d<10.04$.
The value $H_x=0$ follows by differentiating the modular $S$-invariance
of $\vartheta_3(0,\tau)/\eta_D(\tau)$ at $i$. Differentiating
$\eta_D(-1/\tau)=\sqrt{-i\tau}\eta_D(\tau)$ gives $E_2(i)=3/\pi$,
which with $\wp((1-i)/2,i)=0$ gives $H_{h^2}=-\pi/2$.
The positive Eisenstein series for $g_2$, with an explicit tail in the
same script, gives $47.26<g_2(i)/4<47.30$.

Set $\epsilon=3/20000$, $r_h=3/100$, $s=\epsilon/(1/10)=0.0015$,
and $r_2=(r_h/(1/5))^2=0.0225$.
Cauchy's estimates for the even power series of $G$ give, on
$|h|=r_h$, $x=0$,
\[
 |G(h,i)|\geq
 \left(47.26-\frac{500r_2}{1-r_2}\right)r_h^2>0.03217.
\]
For $|x|\leq\epsilon$, the change of $G$ on this circle is at most
$20s/((1-s)(1-r_2))<0.03074$. Rouche's theorem therefore places both
zeros of $G(\cdot,\tau)$ in $|h|<r_h$. They represent the two zeros of
$\wp$ modulo its lattice, since an elliptic function with one double
pole has two zeros with multiplicity.

Write $G=e_2(\tau)+h^2Q(h,\tau)$. The same Cauchy bounds give
\[
 |Q+g_2(i)/4|\leq
 \frac{500(r_2+s/(1-s))}{1-r_2},
\quad
 |e_2(\tau)-idx|\leq
 \frac{2000|x|^2}{1-s}.
\]
At any zero they imply $|x|\leq7|h|^2$ and $|h|^2<1/20000$;
the interval output records bounds $6.124$ and $0.00004434$,
respectively. For $H$, its stated coefficients and Cauchy's estimates
now bound the remainder, together with the change of $1/T$ in the
Gaussian, by $5000|h|^4$. For example a direct upper bound for its
coefficient is
\[
 \frac{0.6(70^2+25\cdot70+25^2)}
 {(1-s)(1-25/20000)}+\frac{7\pi}{1-\epsilon}<5000.
\]
Consequently, for $0<|x|\leq3/20000$,
\[
 \log\mathcal N-\log\sqrt2
 \leq-\frac\pi2|h|^2+5000|h|^4<0.
\]
At $h=0$ the equation $G=0$ forces $x=0$ in this disc by the lower
bound for $e_2$. This proves strict local maximality, with an explicit
radius, and covers the square
$|\operatorname{Re}\tau|\leq10^{-4}$,
$|\operatorname{Im}\tau-1|\leq10^{-4}$.

The grid calculation uses thirty Fourier terms for $\wp$, complex
Newton iteration, and the theta product. Five sample values were
independently recomputed with eighty Fourier terms and the theta series
at ninety-digit precision; discrepancies are below $7\cdot10^{-16}$.
The source, precise grid rule, sampled maximum, and comparisons are in
\texttt{cd\_numerics.py} and its output.
\end{proof}

\begin{theorem}[Certified bound]\label{fam:src:C.2.3}\label{fam:st:C.2.3}\label{ari:C.8.5}\label{push:src:E.4.3}
For the cusp normalization $\beta^{\rm part}+\tau\to0$,
\[
 M=-\frac3\pi\log6=-1.71100425815613411418\ldots,
 \qquad \operatorname*{arg\,max}_{B\setminus\{p_0\}}D_{\rm part}=\{1\}.
\]
The exact parameter discs are
\[
 0<|w|<216,
 \qquad 0<|u|<46656
\]
for the reference linearization $u=w^2$; an arbitrary reference constant
$\lambda_0$ multiplies the latter radius by $|\lambda_0|$.
The proof combines an analytic cusp bound, the quantitative local proof
\ref{fam:src:C.2.2}, and a completed outward-rounded interval cover of the remaining
compact domain.
\end{theorem}
\Uses{the hypotheses, notation and definitions stated here; \cref{fam:src:C.2.2}; \cref{fam:eq:cd-norm}; \cref{fam:src:C.2.1}; \cref{fam:eq:cd-wp}; \cref{fam:src:C.1.1}; \cref{fam:src:C.1.2}}

\begin{proof}
By \eqref{fam:eq:cd-norm} it is enough to prove
$\mathcal N(v,\tau)\leq\sqrt2$ at every zero of $\wp_\tau$, with equality
only in the modular orbit of $i$. Modular transformations, period
translations and sign change preserve this norm. Complex conjugation
also preserves it, through
$(v,\tau)\mapsto(\bar v,-\bar\tau)$, as follows by conjugating the
normally convergent theta and eta series. It therefore suffices to work
on $0\leq\operatorname{Re}\tau\leq1/2$, $|\tau|\geq1$.

For $T\geq2$, the zero-strip argument in \ref{fam:src:C.2.1} improves to
$|\operatorname{Im}v|<b=43/100$. Replace $1/\sinh^2\pi$ in its displayed
error bound by $1/\sinh^2(\pi b)$; the resulting upper bound is less
than $0.316<1/3$, as certified independently in
\texttt{cd\_special\_intervals.py}. Retain the Gaussian in the product
estimate. The function
$\cosh(\pi y)\exp(-\pi y^2/T)$ increases on $0\leq y\leq b$, $T\geq2$:
its logarithmic derivative is $\pi\tanh(\pi y)-2\pi y/T\geq0$, since
$\tanh(\pi y)\geq y$ there. The latter follows from concavity and its
endpoint values. With $r=e^{-4\pi}$ we get
\[
 \mathcal N\leq
 2\cosh(\pi b)e^{-\pi/3-\pi b^2/2}
 \exp\left(\frac{2\cosh(2\pi b)r}{1-r}+
                  \frac{r^2}{1-r^2}\right)
 <1.082<\sqrt2.
\]
The bound decreases for $T\geq2$: the logarithmic derivative of the
exponential $e^{-\pi T/6-\pi b^2/T}$ is negative and the product
majorant decreases with $|q|$. This excludes the whole cusp region,
without a numerical truncation of the base.

We now specify the compact computation as a proof by enclosure.
Its rational starting rectangle is
$\mathcal R=[0,1/2]+i[17/20,2]$, which contains the compact half modular
domain. For every rational rectangle $Q$ retained in this domain, the
program selects a decimal complex centre $v_0$, a decimal nonzero
preconditioner $A$, and a closed complex rectangle $Z$ around $v_0$.
Floating-point Newton iteration only selects these numbers. The proof
uses the subsequent interval tests
\[
 v_0-A\wp(v_0,Q)+(1-A\partial_v\wp(Z,Q))(Z-v_0)
       \subset\operatorname{int}Z,
 \qquad
 \|1-A\partial_v\wp(Z,Q)\|_\infty<1.
\]
Here the norm of multiplication by a complex number is bounded by the
sum of the absolute real and imaginary parts. All displayed sets are
outward-rounded intervals. The mean-value integral along a segment in
$Z$ proves that, for every $\tau\in Q$, the map
$v\mapsto v-A\wp_\tau(v)$ maps $Z$ into itself and is a contraction.
Banach's theorem supplies a genuine zero in $Z$. Both zeros of $\wp$
are related by sign and lattice translations, so their norms agree;
this one enclosure therefore controls the actual period section.

For reproducibility the Fourier remainders are explicit. With
$q=e^{2\pi i\tau}$, $p=q e^{2\pi iv}$, $m=q e^{-2\pi iv}$, let
$r\geq|q|$ and $s\geq\max(|p|,|m|)$ be the interval upper endpoints.
The code requires $s<1$. For $N=24$ put
\[
 S_1(r)=r^{N+1}\left(\frac{N+1}{1-r}+\frac r{(1-r)^2}\right),
\]
\[
 S_2(r)=r^{N+1}\left(\frac{(N+1)^2}{1-r}+
 \frac{2(N+1)r}{(1-r)^2}+\frac{r(1+r)}{(1-r)^3}\right).
\]
The errors in $\wp$ and $\partial_v\wp$ are bounded respectively by
$8\pi^2(S_1(r)+S_1(s))/(1-r)$ and $16\pi^3S_2(s)/(1-r)$.
These follow by absolute summation of \eqref{fam:eq:cd-wp} and its derivative.
They are added as rectangles enclosing the corresponding complex discs.
For the norm, use $h=v-(1-\tau)/2$ and the even theta series
$\vartheta_3=1+2\sum_{n\geq1}e^{\pi i n^2\tau}\cos(2\pi nh)$.
With $T_*\leq\operatorname{Im}\tau$ and $H_*\geq|\operatorname{Im}h|$,
its omitted terms after $n=10$ have absolute value at most
\[
 \frac{2e^{-121\pi T_*+22\pi H_*}}
 {1-e^{-23\pi T_*+2\pi H_*}},
\]
and the denominator is tested positive. The logarithm of the eta
product tail after $n=17$ is bounded by
$r^{18}/((1-r)(1-r^{18}))$. The resulting enclosure of
$e^{-2\pi(\operatorname{Im}h)^2/T}|\vartheta_3/\eta_D|^2$
is required to have upper endpoint strictly below two.

The completed computation \texttt{f\_interval\_parallel.py}, using
\texttt{f\_interval\_cert.py}, runs these tests at thirty decimal digits
and records every rational box. It accepted 19701 norm boxes. A further
36 boxes lie inside the local square of \ref{fam:src:C.2.2} and are covered by that
analytic proof. Another 174 boxes satisfy
$x_{\rm upper}^2+y_{\rm upper}^2<1$ exactly and lie outside the modular
domain. There are zero pending or undecided boxes. The independent
program \texttt{f\_interval\_verify.py} reconstructs the entire binary
partition from its leaves, checks all reported strict interval bounds,
and verifies the exclusions by rational arithmetic. It consumes all
19911 leaves, with exact total area $23/40$ and maximum depth 38.
The largest reported squared-norm upper bound is the exact decimal
rational
\[
 \frac{9999999982782673130994674616307871}
 {5000000000000000000000000000000000}<2.
\]
The JSON certificate, complete outputs, and independently runnable
source programs are included. Thus the compact domain is fully covered;
the exploratory grid is not used as a substitute for any enclosure.
The local proof gives equality only at $\tau=i$, and the cusp and other
compact boxes give strict inequality. This proves the global bound and
its equality set.

Finally \ref{fam:src:C.1.1} evaluates $D$ at $i$ exactly. Since $j(\tau)=1728t$, its
modular orbit maps precisely to $t=1$. Exponentiation gives
$e^{-\pi M}=e^{3\log6}=216$ and $e^{-2\pi M}=6^6=46656$.
The two independent theta calculations and the sixty-digit logarithmic
interval computation in \ref{fam:src:C.1.2} verify the numerical value of $M$ as well
as the exact evaluation; the integer radii are checked against interval
exponentiation in \texttt{cd\_special\_intervals.py}.
\end{proof}


\section{Special values, local expansions and torsion}
\subsection{The normalized Green function and special arguments}
We write $\rho=(1+i\sqrt3)/2$, $T=\operatorname{Im}\tau$, and
\[
 U_\tau(z)=\log N(z,\tau)
 =\log\left|\frac{\vartheta _1(z,\tau)}{\eta(\tau)}\right|
       -\frac{\pi(\operatorname{Im}z)^2}{T}.
\]
Our theta variable has periods $1,\tau$; thus the standard theta series
is evaluated at $\pi z$ when its first argument has period $\pi$.
For a general lattice $\Lambda=\omega(\mathbb Z+\tau\mathbb Z)$ we set
$U_\Lambda(z)=U_\tau(z/\omega)$. The transformation laws of the theta
function show that this definition is independent of the oriented basis.

\begin{proposition}[distribution under isogenies]\label{ari:C.1.1}
Let $\Lambda,\Lambda'\subset\mathbb C$ be lattices and let
$\alpha\ne0$ satisfy $\alpha\Lambda\subset\Lambda'$. Write
$\phi:\mathbb C/\Lambda\to\mathbb C/\Lambda'$ for $z\mapsto\alpha z$,
$K=\alpha^{-1}\Lambda'/\Lambda$, and
$d=[\Lambda':\alpha\Lambda]$. Then
\begin{equation}\label{ari:eq:cs-isogeny}
 U_{\Lambda'}(\alpha z)=\sum_{v\in K}U_\Lambda(z+v).
\end{equation}
The equality holds away from the zeros and as an equality of locally
integrable functions at them. If $\Delta_\Lambda=g_2(\Lambda)^3-
27g_3(\Lambda)^2$, then
\[
 \prod_{0\ne v\in K}N_\Lambda(v)
 =|\alpha|\left|\frac{\Delta_{\Lambda'}}{\Delta_\Lambda}\right|^{1/12}.
\]
In particular, for every nonzero endomorphism of one elliptic curve,
$\prod_{v\in\ker\alpha,\ v\ne0}N(v)=|\alpha|$.
\end{proposition}
\Uses{the hypotheses, notation and definitions stated here; the explicit constructions and calculations in the proof}

\begin{proof}
The Jacobi product, with $q=e^{2\pi i\tau}$ and
$X=e^{2\pi iz}$, gives
\[
 \vartheta _1(z,\tau)/\eta(\tau)
 =i q^{1/12}X^{-1/2}(1-X)
     \prod_{m\ge1}(1-q^mX)(1-q^mX^{-1}),
\]
where a harmless fixed unit changes with the theta convention. For
$0<y<T$, integration of the logarithm over $z=x+iy$, $0\le x\le1$,
gives $-\pi T/6+\pi y-\pi y^2/T$. Each product factor has mean logarithm
zero by the absolutely convergent logarithmic series; the endpoint
singularities are integrable. Integrating in $y$ proves
$\int U_\Lambda\,d\mu_\Lambda=0$ for Haar probability measure.
Moreover
\[
 \Delta_{\rm Eucl}U_\Lambda
   =2\pi\delta_0-\frac{2\pi}{\operatorname{area}\Lambda},\qquad
 U_\Lambda(z)=\log|z|+\frac1{12}\log|\Delta_\Lambda|+o(1).
\]
The second formula follows from
$\vartheta _1'(0,\tau)=2\pi\eta(\tau)^3$ and
$\Delta_{\mathbb Z+\tau\mathbb Z}=(2\pi)^{12}\eta(\tau)^{24}$,
then scaling by $\omega$. The two sides of
\eqref{ari:eq:cs-isogeny} have the same logarithmic masses at $K$ and the
same constant Laplacian, since
$d=|\alpha|^2\operatorname{area}\Lambda/
\operatorname{area}\Lambda'$. Their difference is a smooth harmonic
function on the compact torus, hence constant. Both means vanish:
translations preserve Haar measure and an isogeny sends Haar
probability measure to Haar probability measure. The constant is zero.
Taking $z\to0$ and cancelling the logarithmic singularity gives the
product. This also checks the precise normalization of the special-value calculation in Proposition~\ref{push:src:D.4.5};
its kernel, Laplacian and mean-zero argument are all valid.
\end{proof}
\begin{proposition}[the phases and the two periods]\label{ari:C.1.2}
With the branch continued from the specified cusp, the exact values are
\begin{align*}
 \frac{\vartheta _1((2-\rho)/3,\rho)}{\eta(\rho)}
   &=e^{\pi i/12}\,3^{1/4}e^{\pi/(6\sqrt3)},\\
 \frac{\vartheta _1((1-i)/2,i)}{\eta(i)}
   &=\sqrt2\,e^{\pi/4},\\
 \beta^{\rm part}(z_1)
   &=-\frac12+i\left(\frac1{\sqrt3}-\frac3{2\pi}\log48\right),\\
 \beta^{\rm part}(z_2)
   &=i\left(\frac32-\frac3\pi\log6\right).
\end{align*}
Consequently $N(z_1)=3^{1/4}$, $N(z_2)=\sqrt2$,
$D_{\rm part}(z_1)=-3\log48/(2\pi)$, and
$D_{\rm part}(z_2)=-3\log6/\pi$.
\end{proposition}
\Uses{the hypotheses, notation and definitions stated here; \cref{ari:C.1.1}}

\begin{proof}
The special-value calculation in Proposition~\ref{push:src:D.4.5}
uses multiplication by $1+\rho$ has matrix
$\left(\begin{smallmatrix}1&-1\\1&2\end{smallmatrix}\right)$,
its kernel is $\{0,\mu,-\mu\}$, and $(1+\rho)\mu=1$ for
$\mu=(2-\rho)/3$. Statement \ref{ari:C.1.1} gives $N(\mu)^2=\sqrt3$.
For the requested parallel calculation at $i$, multiplication by $1+i$
has kernel $\{0,(1-i)/2\}$ and absolute value $\sqrt2$; the same
identity gives $N((1-i)/2,i)=\sqrt2$ without a separate theta-constant
evaluation.

At $\rho$, $\mu=1/2-i/(2\sqrt3)$ and $q=-e^{-\pi\sqrt3}$. In
\[
 R(z,\tau)=2e^{\pi i\tau/6}\sin(\pi z)
 \prod_{m\ge1}\bigl(1-2q^m\cos(2\pi z)+q^{2m}\bigr),
\]
$\sin(\pi\mu)$ is positive real. The factors are
$(1+q^m e^{\pi/\sqrt3})(1+q^m e^{-\pi/\sqrt3})$, both positive:
when $m$ is odd the larger subtracted number is at most
$e^{-2\pi/\sqrt3}<1$. Thus the phase is exactly $\pi/12$ modulo
$2\pi$. There is no winding ambiguity in the continued logarithm.
Indeed continue from the cusp along $\tau=1/2+iT$,
$T\ge\sqrt3/2$. Real conjugation of the rhombic lattice keeps the
chosen zero on $\mu=1/2-iy$, $0<y<T$. It cannot leave that component:
a collision of the two zeros would require $g_3=0$, which does not
occur on this ray. Indeed, writing $q=-r$, $0<r<1/200$,
the tail after the first term of the $E_6$ series has absolute value
at most $504\cdot32r^2/((1-r^2)(1-8r))<1/2$, while its
first term is positive. Thus $E_6>1/2$. Hitting $y=0$ would again give a half-period zero, and
hitting $y=T$ would give another half-period zero. Either event
would force $g_3=0$. Each displayed product factor remains
positive because $e^{-2\pi T}e^{2\pi y}<1$. The cusp phase tends to
$\pi/12$, so the continued argument has that value. On the imaginary
ray $\tau=iT$, $T\ge1$, the same branch is
$\mu=1/2-iy$ with $0<y\le T/2$; every product factor is positive,
so its argument at $i$ is zero. These checks supply the branch step
used in R3. The Gaussian factors are respectively
$e^{\pi/(6\sqrt3)}$ and $e^{\pi/4}$. Substitution in the definition of
$\beta^{\rm part}$ proves all four identities.
\end{proof}

\begin{proposition}[numerical cross-check]\label{ari:C.1.3}
The independent analytic computation in
\path{code/Cspecial_torsion_numerics.py} gives, to fifty decimals,
\begin{align*}
 N(z_1)&=1.31607401295249246081921890179699905516006859020582\ldots,\\
 N(z_2)&=1.41421356237309504880168872420969807856967187537695\ldots,\\
 \operatorname{Im}\beta^{\rm part}(z_1)
  &=-1.27101206057537346286579673538412682778809507094949\ldots,\\
 \operatorname{Im}\beta^{\rm part}(z_2)
  &=-0.21100425815613411418178314271018934802718590587233\ldots.
\end{align*}
The theta/eta evaluations and exact expressions differ by less than
$10^{-100}$ at working precision $110$ decimal digits.
\end{proposition}
\Uses{the hypotheses, notation and definitions stated here; \cref{ari:C.1.2}}

\begin{proof}
The program directly evaluates the theta series and the eta product,
then compares the displayed closed forms. The eta product has $160$
factors. Since $|q|\le e^{-\pi\sqrt3}<1/200$, the omitted logarithm is
at most $\sum_{m>160}|q|^m/(1-|q|^m)<10^{-368}$.
The output records the measured floating-point discrepancies. The
exact values are established by \ref{ari:C.1.2}; this computation is their
independent numerical check, not an additional rounding certificate.
\end{proof}


\subsection{Division identities and the entire torsion locus}
We use throughout this subsection the rational model
\[
 E_t:\quad y^2=x^3+A(t)x+B(t),\qquad
 A=-\tfrac34t^3(t-1),\quad B=\tfrac14t^4(t-1)^2,
 \quad P=(0,t^2(t-1)/2).
\]
Its discriminant is $\Delta(t)=27t^8(t-1)^3$.
The change of complex Weierstrass model does not change $N$.


\paragraph{Cited input.} \label{ari:C.4.1}The source statement and its expanded derivation are retained in the proof compendium, Paper~I, label \nolinkurl{ari:C.4.1}.\par

\Uses{the hypotheses, notation and definitions stated here; \citet[Remark~2.2 and Proposition~2.3]{Shin2015}; \cref{app:rederive:ari-C.4.1}}


The Siegel-function identities are those of \citet[Remark~2.2 and Proposition~2.3]{Shin2015}, with the local-height normalization proved explicitly. A complete re-derivation is given in Appendix~\ref{app:rederive:ari-C.4.1}.


\begin{proposition}[the full division identity]\label{ari:C.4.2}
Let $\psi_m$ be the standard division polynomial of a smooth
Weierstrass model, normalized by $\psi_1=1$, $\psi_2=2y$.
For $m\ge2$ and $Q,mQ\ne O$,
\begin{equation}\label{ari:eq:cs-division}
 N(Q)^{m^2}
       =N(mQ)\frac{|\Delta_E|^{(m^2-1)/12}}{|\psi_m(Q)|}.
\end{equation}
For the displayed family this specializes to
\[
 N(P)^4=3^{3/4}|t-1|^{-1/4}N(2P),\qquad
 N(P)^9=16|t|^{-2/3}N(3P).
\]
If $P_t$ has exact order $n$ and $m\ge2$ satisfies
$m\equiv1$ or $-1\pmod n$, then
\begin{equation}\label{ari:eq:cs-torsionvalue}
 N(P_t)^{12}
       =\frac{|\Delta(t)|}{|\psi_m(P_t)|^{12/(m^2-1)}}.
\end{equation}
The positive value on the right is independent of $m$.
\end{proposition}
\Uses{the hypotheses, notation and definitions stated here; \cref{ari:C.1.1}}

\begin{proof}
Use the complex uniformization $x=\wp_\Lambda(z)$,
$y=\wp'_\Lambda(z)/2$. The rational function $\psi_m$ has divisor
$\sum_{0\ne Q\in E[m]}(Q)-(m^2-1)(O)$ and leading Laurent term
$(-1)^{m-1}m z^{1-m^2}$. One way to verify both normalizations without
a model convention is to form
$(-1)^{m-1}\sigma_\Lambda(mz)/\sigma_\Lambda(z)^{m^2}$.
Its lattice factors cancel, its zeros are exactly the nonzero
$m$-division points, and its expansion at zero is the displayed
one. The addition formula for $\sigma$ gives the standard
division-polynomial recurrence with initial values $1,2y$, so it is
the stated rational function.

The function
\[
 m^2U_\Lambda(z)-U_\Lambda(mz)
                    +\log|\psi_m(Q(z))|
\]
is harmonic after its logarithmic singularities cancel. Its
Laplacian has zero constant term as well, so it is constant.
The expansion in \ref{ari:C.1.1} and the leading term of $\psi_m$ make that
constant $(m^2-1)\log|\Delta_E|/12$. This proves
\eqref{ari:eq:cs-division}. Direct substitution of
\[
 \psi_2(P)=t^2(t-1),\qquad
 \psi_3(P)=-9t^6(t-1)^2/16
\]
gives the two special identities. For $m\equiv\pm1\pmod n$,
$mP=\pm P$ and $N$ is even. Cancelling $N(P)$ in
\eqref{ari:eq:cs-division} gives \eqref{ari:eq:cs-torsionvalue}.
Every right side is the same positive number $N(P)^{12}$, which
proves independence of $m$, not merely numerical agreement.
\end{proof}

\begin{proposition}[algebraicity and explicit fields for every order]\label{ari:C.4.3}
Let $t_0\in\mathbb C\setminus\{0,1\}$ and suppose $P_{t_0}$ has
exact order $n\ge2$. Then $t_0$ and $N(P_{t_0})$ are algebraic.
Put $K_0=\mathbb Q(t_0,\bar t_0)$,
\[
 W=\Delta(t_0)\Delta(\bar t_0)>0,\qquad
 V=\psi_{n+1}(P_{t_0})\psi_{n+1}(P_{\bar t_0})>0.
\]
With positive real radicals,
\[
 N(P_{t_0})^{12n}=W^{n/2}V^{-6/(n+2)}
 \ \in\ K_0\bigl(W^{1/2},V^{1/(n+2)}\bigr).
\]
This supplies an explicit field, without a claim that it is minimal.
More generally, if $m\equiv\pm1\pmod n$, then
\[
 \lambda_\infty(P_{t_0})
 =\frac{\log|\psi_m(P_{t_0})|}{m^2-1}
       -\frac1{12}\log|\Delta(t_0)|.
\]
These formulas determine the height and $N^{12}$ on the entire
torsion locus.
\end{proposition}
\Uses{the hypotheses, notation and definitions stated here; \cref{ari:eq:cs-torsionvalue}}

\begin{proof}
The division polynomial $\psi_n(P_t)$ is a nonzero rational
polynomial in $t$. To check nonvanishing for every $n$, a hypothetical
torsion identity would express the elliptic logarithm as
$a\tau+b$ with constant rational $a,b$. Its bounded cusp limit
forces $a=0$, but that limit has nonzero imaginary part, a
contradiction. Thus $t_0$ is a zero of a nonzero polynomial over
$\mathbb Q$. Formula \eqref{ari:eq:cs-torsionvalue} shows algebraicity
of $N$ using only positive rational powers of nonzero algebraic
absolute values. Taking $m=n+1$, raising to the $n$th power, and
using $|\Delta|=W^{1/2}$, $|\psi_{n+1}|=V^{1/2}$ gives the exact
field statement. The local-height formula is its logarithm.

Here is a polynomial relation involving only $t$ and $r=N^{12}$,
if one prefers that form to positive radicals. Remove the factors
at $0,1$ from $\psi_n(P_t)$ and divide successively by the
corresponding polynomials for proper divisors of $n$; let $F_n$ be
the resulting primitive polynomial. At a point of exact order
$d\mid n$, the formal group gives
$z([n/d]Q)=(n/d)z(Q)+O(z(Q)^2)$, so the vanishing multiplicities for
$\psi_d$ and $\psi_n$ agree. Hence this cancellation leaves exactly
the exact-order-$n$ divisor, with its multiplicities. Set
$m=n+1$ and $e=m^2-1$. Clearing rational denominators in
\[
 \operatorname{Res}_S\left(
 F_n(S),\
 r^{2e}\{\psi_m(P_T)\psi_m(P_S)\}^{12}
              -\{\Delta(T)\Delta(S)\}^{e}\right)
\]
gives a polynomial in $\mathbb Q[T,r]$ that vanishes at
$(t_0,N(P_{t_0})^{12})$, since $S=\bar t_0$ is a common root.
It is nonzero: at an exact-order-$n$ root no $\psi_m$ vanishes,
so the leading coefficient in $r$ of the resultant is nonzero.
This is an explicit recurrence-and-resultant prescription for
every $n$, including nonreal torsion parameters.
\end{proof}

\begin{proposition}[the five values and all orders through eight]\label{ari:C.4.4}
The exact special values are
\[
 \begin{array}{c|c|c}
 t&\operatorname{ord}(P)&N(P)^{12}\\ \hline
 0&3&27\\
 1&2&64\\
 32/5&4&10\\
 64(47-27\sqrt3)/11&6&17+10\sqrt3\\
 64(47+27\sqrt3)/11&6&10\sqrt3-17 .
 \end{array}
\]
At $t=0,1$ the orders refer to the smooth models after the specified
base changes. On the smooth base, $F_2=F_3=1$ and
\[
 \begin{aligned}
 F_4&=5t-32,\\
 F_5&=89t^2+4736t-4096,\\
 F_6&=11t^2-6016t+8192,\\
 F_7&=881t^4+4745344t^3-12079104t^2\\
    &\hspace{13mm}+24641536t-16777216,\\
 F_8&=49t^4-2614336t^3-7864320t^2\\
    &\hspace{13mm}+19398656t-8388608.
 \end{aligned}
\]
All these polynomials are squarefree. At each root of exact order
$n=4,5,6,7,8$, the positive number $r=N^{12}$ is respectively
determined by
\[
 \begin{array}{c|l}
 n&\text{equation for }r\\ \hline
 4&r=64/|t|\\
 5&r^5=432^5/\bigl((1-t)(5t-32)^4\bigr)\\
 6&r^2=1728^2/F_5(t)\\
 7&r^{35}=2^{192}3^{81}/
       \bigl(|t|^8|t-1|^3|F_6(t)|^{12}\bigr)\\
 8&r^4=1728^4/|F_7(t)|.
 \end{array}
\]
The order-five roots are real and below one; the order-six roots
are real and above one. The two complex order-seven roots are
included in the formula.
\end{proposition}
\Uses{the hypotheses, notation and definitions stated here; \cref{ari:eq:cs-torsionvalue}; \cref{ari:C.1.2}; \cref{ari:C.4.3}}

\begin{proof}
The division-polynomial recurrence, evaluated at $x=0$,
$y=t^2(t-1)/2$, gives
\[
 \begin{aligned}
 \psi_4(P)&=-t^{10}(t-1)^4F_4/32,\\
 \psi_5(P)&=t^{16}(t-1)^6F_5/4096,\\
 \psi_6(P)&=9t^{24}(t-1)^9F_6/65536,\\
 \psi_7(P)&=-t^{32}(t-1)^{12}F_7/16777216,\\
 \psi_8(P)&=-t^{42}(t-1)^{16}F_4F_8/268435456.
 \end{aligned}
\]
The recursion and all coefficients are checked over $\mathbb Q[t]$
by \path{Cspecial_division.py}. The same program computes
$\gcd(F_n,F_n')=1$ and excludes common factors for proper orders.
Thus no point or multiplicity is omitted. Applying
\eqref{ari:eq:cs-torsionvalue} with $m=n-1$ gives every displayed
formula. The signs for $F_5,F_6$ follow directly from their quadratic
formula; in particular $F_5$ is positive at the two roots of $F_6$.
At order six, substituting
$t=64(47\mp27\sqrt3)/11$ gives the exact identities
\[
 F_5(t)(17\pm10\sqrt3)^2=1728^2.
\]
The positive roots are $17+10\sqrt3$ and $10\sqrt3-17$; their
product is $11$. The substitution is checked in
\path{Cspecial_relations.py}. The order-four value follows from
$64/(32/5)=10$, and the two central values are \ref{ari:C.1.2}.

For clarity, the recurrence used by the exact program is
\[
 \begin{aligned}
 \psi_{2m+1}&=\psi_{m+2}\psi_m^3-\psi_{m-1}\psi_{m+1}^3,\\
 \psi_{2m}&=\frac{\psi_m}{\psi_2}
   (\psi_{m+2}\psi_{m-1}^2-\psi_{m-2}\psi_{m+1}^2).
 \end{aligned}
\]
Its initial values at $P$ are $\psi_0=0,\psi_1=1,\psi_2=2y$,
$\psi_3=-A^2$, $\psi_4=4y(-8B^2-A^3)$. No approximate algebra
enters this calculation. The same code records $\psi_9$, and the
separate file \texttt{Cspecial\_psi16.out} records the factored
$\psi_{16}$ for the duplication ladder. A rational polynomial
relation for the nonreal order-seven values is explicitly computed
by \path{Cspecial_relations.py}, by the resultant prescription
of \ref{ari:C.4.3}.
\end{proof}

\begin{proposition}[complete numerical verification through order eight]\label{ari:C.4.5}
All $13$ smooth exact-order roots for $2\le n\le8$ were evaluated
at $110$ decimal digits independently of the algebraic norm formula.
The discrepancy in $N^{12}$ is below $10^{-90}$ at every root;
the period coordinates also verify $nP=O$ to that tolerance.
There are $0,0,1,2,2,4,4$ roots for $n=2,\ldots,8$.
\end{proposition}
\Uses{the hypotheses, notation and definitions stated here; the explicit constructions and calculations in the proof}

\begin{proof}
The script \path{Cspecial_torsion_numerics.py} computes the
roots of $F_n$ to $115$ digits and constructs periods of
$y^2=x^3+Ax+B$ using its three roots $e_1,e_2,e_3$. If
$k=(e_2-e_3)/(e_1-e_3)$ and $K$ denotes the complete elliptic
integral with parameter $k$, it uses
\[
 \tau=iK(1-k)/K(k),\qquad
 v=\frac{F\bigl(\arcsin\sqrt{(e_1-e_3)/(-e_3)},k\bigr)}{2K(k)}.
\]
The identity
$\wp(u)=e_3+(e_1-e_3)/\operatorname{sn}^2
(\sqrt{e_1-e_3}\,u,k)$ verifies that this is a zero of $\wp$.
Permutations of the roots are tried until $\operatorname{Im}\tau>0$;
modular reduction and lattice translations improve numerical
conditioning and leave $N$ unchanged. The program then evaluates
the theta quotient directly. It checks both the norm formula and
that the coefficients of $nv$ in the basis $(\tau,1)$ are integral.
The full output gives at least $70$ digits for every $t,\tau,v,N$
and $N^{12}$, including both nonreal conjugates. The largest measured
norm discrepancy is below $2\cdot10^{-108}$. These are numerical
checks of the exact theorem, not numerical proofs of it.
\end{proof}


\section{An analytic proof of the elliptic maximum}

Throughout this section $q=e^{2\pi i\tau}$, $T=\Im\tau$,
$\Delta_0=q\prod_{n\ge1}(1-q^n)^{24}$, and
$t=E_4^3/(1728\Delta_0)$. Thus $t-1=E_6^2/(1728\Delta_0)$.
The discriminant $\Delta_0$ is the modular discriminant and must not
be confused with the discriminant $27t^8(t-1)^3$ of the rational
Weierstrass equation. We use the division identity
\eqref{ari:eq:cs-division}, already proved above, and the modular
invariance of $N$. No interval cover or numerical root computation
enters the following proof.

\begin{proposition}[three uniform bounds on the elliptic norm]\label{ari:C.8.1}
For every $z\in\C$ and every $\tau$ in the standard modular
fundamental region, we have
\[
 \begin{array}{c|ccc}
 \text{restriction}&T\le2&T\le7/5&T\le11/10\\ \hline
 N(z,\tau)&<17/10&<3/2&<143/100.
 \end{array}
\]
If $\wp_\tau(z)=0$ and $T\ge2$, then $N(z,\tau)<541/500$.
\end{proposition}
\Uses{the hypotheses, notation and definitions stated here; \cref{ari:eq:ca-sup}; \cref{ari:eq:ca-exp}}

\begin{proof}
Taking absolute values in the theta series and completing the square
gives
\[
 e^{-\pi(\Im z)^2/T}|\vartheta_1(z,\tau)|
 \le\sum_{n\in\Z}
 e^{-\pi T(n+1/2+\Im z/T)^2}
 \le\sum_{n\in\Z}e^{-\pi Tn^2}.
\]
The last inequality follows from Poisson summation: the Fourier
coefficients of the periodized Gaussian are positive, so its maximum
is at an integer. The eta product and $n^2\ge1+3(n-1)$ for $n\ge1$
therefore imply, on $a\le T\le b$,
\begin{equation}\label{ari:eq:ca-sup}
 N(z,\tau)\le
 \exp\!\left(\frac{r_a}{(1-r_a)^2}\right)
 \max\{F(a),F(b)\},\qquad
 F(T)=e^{\pi T/12}
 \left(1+\frac{2e^{-\pi T}}{1-e^{-3\pi T}}\right),
\end{equation}
where $r_a=e^{-2\pi a}$. Indeed,
$-\sum\log(1-r^n)\le r/(1-r)^2$. The function $F$ is convex,
as follows by expanding its denominator as a geometric series;
each summand is an exponential of a real linear function of $T$.
We use $6/7<\sqrt3/2$ and the following rational bounds:
\[
\begin{array}{c|@{\,}c@{\,}c@{\,}c@{\,}c@{\,}c}
T&6/7&1&11/10&7/5&2\\ \hline
F(T)&<14211/10000&<14116/10000&<14180/10000&<14782/10000&<16944/10000.
\end{array}
\]
The eta factors in \eqref{ari:eq:ca-sup} are less than $10047/10000$
for $a=6/7$ and less than $10019/10000$ for $a=1$.
Multiplying these fractions proves the asserted bounds, using
$[6/7,1]$ and $[1,2]$ separately for the first one.

We specify elementary arithmetic for all the scalar estimates in
this proof. One may use $314159/100000<\pi<355/113$ and, for
$0\le x\le20$, the Taylor bounds
\begin{equation}\label{ari:eq:ca-exp}
 S_{80}(x)\le e^x\le
 S_{80}(x)+\frac{x^{81}}{81!}\frac1{1-x/82},
 \qquad S_{80}(x)=\sum_{j=0}^{80}\frac{x^j}{j!}.
\end{equation}
Taking reciprocals bounds negative exponentials. The remainder
bound follows by comparison of consecutive terms, not by numerical
sampling. These inequalities verify each displayed rational bound;
much shorter Taylor sums also suffice. Thus the decimal diagnostics
in the accompanying programs are dispensable.

For the cusp we check directly the analytic step of supplied K6.
Choose $|\Im z|\le T/2$ by translation, put $b_0=43/100$ and
$r_0=e^{-4\pi}$. The Fourier series
\[
 \frac{\wp_\tau(z)}{\pi^2}
 =\csc^2(\pi z)-\frac13
   +8\sum_{n\ge1}\frac{nq^n}{1-q^n}
          \bigl(1-\cos(2\pi nz)\bigr)
\]
shows, if $|\Im z|\ge b_0$, that
\begin{align*}
 \left|\frac{\wp_\tau(z)}{\pi^2}+\frac13\right|
 &\le \frac1{\sinh^2(\pi b_0)}
 +8\left\{\frac{r_0}{(1-r_0)^3}
 +\frac1{2(1-r_0)}
 \left(\frac{e^{-2\pi}}{(1-e^{-2\pi})^2}
       +\frac{e^{-6\pi}}{(1-e^{-6\pi})^2}\right)\right\}
 <\frac13.
\end{align*}
The two terms in the last parentheses bound respectively
$r^n e^{2\pi n|\Im z|}$ and $r^n e^{-2\pi n|\Im z|}$.
This inequality excludes a zero. Hence a zero has
$|\Im z|<b_0$. The Jacobi product, $\log(1+x)\le x$, and
the monotonicity of $\cosh(\pi y)e^{-\pi y^2/T}$ on
$0\le y\le b_0$ give
\[
 N(z,\tau)\le
 2\cosh(\pi b_0)e^{-\pi/3-\pi b_0^2/2}
 \exp\!\left(
 \frac{2\cosh(2\pi b_0)r_0}{1-r_0}
 +\frac{r_0^2}{1-r_0^2}\right)<\frac{541}{500}.
\]
For the stated monotonicity, the concave function
$\tanh(\pi y)-y$ is zero at zero and positive at $b_0$.
For $T\ge2$ the right side before substituting $T=2$ decreases:
the derivative of $-\pi T/6-\pi b_0^2/T$ is negative and
both product errors decrease. The two strict scalar inequalities
follow from \eqref{ari:eq:ca-exp}. This proves the cusp assertion
without using any certificate from K6.
\end{proof}

\begin{proposition}[an explicit local maximum on a larger disc]\label{ari:C.8.2}
If $|\tau-i|\le1/50$ and $\wp_\tau(z)=0$, then
$N(z,\tau)\le\sqrt2$, with equality exactly when $\tau=i$.
\end{proposition}
\Uses{the hypotheses, notation and definitions stated here; \cref{ari:eq:ca-exp}; \cref{ari:eq:ca-derivatives}; \cref{ari:eq:ca-GHbounds}; \cref{ari:eq:ca-H}; \cref{ari:eq:ca-delta-h}}

\begin{proof}
Write $\delta=\tau-i$, $h=z-(1-\tau)/2$ and
\[
 G(h,\tau)=\wp_\tau((1-\tau)/2+h),\qquad
 H(h,\tau)=\log\frac{\vartheta_3(h,\tau)}{\eta(\tau)}.
\]
We choose the holomorphic logarithm whose value at $(0,i)$ is
$\log\sqrt2$. The half-period theta identity gives
\begin{equation}\label{ari:eq:ca-H}
 \log N(z,\tau)=\Re H(h,\tau)-\frac{\pi(\Im h)^2}{T}.
\end{equation}
Both $G$ and $H$ are even in $h$. On the polydisc
$|h|\le1/4$, $|\delta|\le1/10$, they satisfy
\begin{equation}\label{ari:eq:ca-GHbounds}
 |G|<30,\qquad |H|<2/3.
\end{equation}
Here are sufficient series estimates, including their margins.
Put
\[
 A_k=\sum_{n\ge1}n^k e^{-9\pi n^2/10}\cosh(\pi n/2),
 \quad a_0=2A_0,\quad a_1=4\pi A_1,\quad a_2=8\pi^2A_2.
\]
Then $a_0<3/10$, $a_1<19/10$, $a_2<12$.
The first two summands suffice to verify these inequalities: for
$n\ge3$, each of the two exponentials in the hyperbolic cosine
has successive-term ratio at most $e^{-29\pi/5}$, and the
factor $n^k$ increases this ratio by at most $(4/3)^2$.
We use \eqref{ari:eq:ca-exp} for the first terms and this geometric
bound for the tail. With $r=e^{-9\pi/5}$, the identities
\[
 G=-\partial_h^2\log\vartheta_3-\frac{\pi^2}{3}E_2,
 \qquad
 |E_2|\le1+\frac{24r}{(1-r)^3}
\]
give
\[
 |G|\le\frac{a_2}{1-a_0}+\frac{a_1^2}{(1-a_0)^2}
       +\frac{\pi^2}{3}\left(1+\frac{24r}{(1-r)^3}\right)<30.
\]
Also
\[
 |H|\le-\log(1-a_0)+\frac{11\pi}{120}
                    +\frac{r}{(1-r)^2}<\frac23.
\]
The logarithm is well defined since $|\vartheta_3-1|<3/10$;
the displayed eta product fixes its branch. A geometric bound
for $-\log(1-x)=\sum_{n\ge1}x^n/n$ verifies the last estimate.

Let
\[
 e(\tau)=\frac{\pi^2}{3}(\vartheta_2(0,\tau)^4-
                              \vartheta_4(0,\tau)^4),\qquad
 E(\tau)=\frac{\pi^2}{3}E_2(\tau),\qquad
 g=\frac{g_2(i)}4.
\]
On $|\delta|\le1/50$ the following loose bounds hold:
\begin{equation}\label{ari:eq:ca-derivatives}
 |e'|<14,\quad |e''|<60,\quad |g_2'|<500,
 \quad |E'|<2,\quad |E_6''|<100,
 \qquad 47<g<95/2.
\end{equation}
For completeness these are scalar Fourier-series bounds, not
interval computations. At $T\ge49/50$, upper bounds for
$|\vartheta_j|,|\vartheta_j'|,|\vartheta_j''|$ are
\[
\begin{array}{c|ccc}
j&0\text{ derivatives}&1\text{ derivative}&2\text{ derivatives}\\ \hline
2&929/1000&742/1000&670/1000\\
4&1093/1000&290/1000&911/1000.
\end{array}
\]
They follow by differentiating the absolutely convergent theta
series, taking absolute values, retaining two terms, and using
the same successive-term bound. Differentiating the fourth
powers proves the first two bounds in \eqref{ari:eq:ca-derivatives}.
For $r\le1/470$, the remaining derivative majorants are
\begin{align*}
 |g_2'|&\le640\pi^5\frac{r(1+11r+11r^2+r^3)}{(1-r)^7}<500,\\
 |E'|&\le16\pi^3\frac{r(1+r)}{(1-r)^5}<2,\\
 |E_6''|&\le2016\pi^2
 \frac{r(1+r)(1+120r+1191r^2+2416r^3+1191r^4+120r^5+r^6)}
 {(1-r)^{11}}<100.
\end{align*}
These follow by summing $\sum n^k r^n$ and bounding
$(1-r^n)^{-1}$ by $(1-r)^{-1}$. Substitution of $r=1/470$
and $\pi<355/113$ verifies the rational inequalities.
Finally $g=\pi^4E_4(i)/3$ and
$E_4=1+240\sum n^3q^n/(1-q^n)$ give $47<g<95/2$.
The half-period identity used in Lemma~\ref{push:src:D.2.1},
obtained by differentiating $e$ and using the theta heat equation, reads
\[
 e(i)=0,\qquad e'(i)=id,\qquad d=\frac{2g}{3\pi}>99/10.
\]

The differential equation $\wp''=6\wp^2-g_2/2$ determines the
first coefficients in the even expansion:
\[
 G(h,\tau)=e+a h^2+ea h^4+\mathcal R_G,
 \qquad a=3e^2-g_2/4.
\]
Consequently, on $|\delta|\le1/50$,
\[
 |e|\le7/25,\quad |a+g|\le\frac52+3(7/25)^2=\frac{3419}{1250},
 \quad |a|\le\frac{62794}{1250}.
\]
Cauchy's estimate from \eqref{ari:eq:ca-GHbounds} yields
\[
 |\mathcal R_G|\le
 \frac{30(4|h|)^6}{1-16|h|^2}.
\]
On $|h|=1/10$ the term $ah^2$ has modulus at least
$55331/125000>0.4426$, whereas the sum of the other terms is at most
\[
 \frac7{25}+\frac7{25}\frac{62794}{1250}\frac1{10000}
       +\frac{30(4/25)^3}{1-4/25}<0.4277.
\]
Rouch\'e's theorem puts two zeros, with multiplicity, in
$|h|<1/10$. The disc injects into the torus and $\wp$ has exactly
two zeros on the torus, so it contains both of them.

Write $G=e+h^2Q$. The preceding estimates give, for $|h|\le1/10$,
\[
 |Q+g|\le
 \frac{3419}{1250}+\frac7{25}\frac{62794}{1250}\frac1{100}
       +\frac{30(4/25)^3}{(1-4/25)/100}<\frac{88}{5}.
\]
Taylor's formula and $|e''|<60$ give
$|e|\ge(99/10-3/5)|\delta|=(93/10)|\delta|$.
At either zero, therefore,
\begin{equation}\label{ari:eq:ca-delta-h}
 |\delta|<7|h|^2,
\end{equation}
because $(95/2+88/5)/(93/10)=7$.

Expand
$H=H_0(\tau)+H_2(\tau)h^2+H_4(\tau)h^4+\mathcal R_H$.
The heat equation and the identity for $G$ give
\[
 H_0'=\frac{ie}{4\pi},\quad
 H_2=-\frac{e+E}{2},\quad H_4=-\frac a{12}.
\]
Thus $H_0(i)=\log\sqrt2$, $H_0'(i)=0$,
$|H_0''|<6/5$, $H_2(i)=-\pi/2$, $|H_2'|<8$, and
$|H_4|<21/5$. Cauchy's estimate gives
\[
 |\mathcal R_H|\le
 \frac{(2/3)(4|h|)^6}{1-16|h|^2}
 \le\frac{2048}{63}|h|^4
 \quad (|h|\le1/10).
\]
Combining \eqref{ari:eq:ca-H} and \eqref{ari:eq:ca-delta-h}, we obtain
\begin{align*}
 \log N-\log\sqrt2
 &\le-\frac\pi2\Re(h^2)-\pi(\Im h)^2+C|h|^4\\
 &=-\frac\pi2|h|^2+C|h|^4,\\
 C&=\frac{147}{5}+56+\frac{21}{5}
                 +\frac{2048}{63}+\frac{1100}{49}<150.
\end{align*}
The final term of $C$ bounds the change of $1/T$ from $1$:
$7\pi/(1-1/50)<1100/49$. Since $150|h|^2<3/2<\pi/2$,
the result is strict unless $h=0$. Equation
\eqref{ari:eq:ca-delta-h} then gives $\delta=0$, where equality
holds by the checked isogeny computation at $i$.
\end{proof}

\begin{proposition}[the comparison of the two local discs]\label{ari:C.8.3}
Every complex number $t_0$ with $|t_0-1|\le1/500$ has a
representative $\tau$ with $|\tau-i|<1/50$ and $j(\tau)=1728t_0$.
Consequently every zero of $\wp$ for that modulus satisfies
$N\le\sqrt2$, with equality only for $t_0=1$.
\end{proposition}
\Uses{the hypotheses, notation and definitions stated here; \cref{ari:eq:ca-derivatives}; \cref{ari:C.8.2}}

\begin{proof}
At $i$, $E_6=0$ and the Ramanujan differential identity gives
$|E_6'(i)|=\pi E_4(i)^2>13/2$, since $E_4(i)>29/20$.
On $|\delta|=1/50$, \eqref{ari:eq:ca-derivatives} gives
\[
 |E_6(\tau)|>\frac{13}{100}-\frac1{50}=\frac{11}{100}.
\]
On the whole disc, $|q|<1/470$ and
$|\Delta_0|\le |q|\exp(24|q|/(1-|q|))<1/400$.
It follows on the boundary that
\[
 |t-1|=\frac{|E_6|^2}{1728|\Delta_0|}
 >\frac{121}{43200}>\frac1{500}.
\]
Rouch\'e applied to $E_6'(i)\delta$ shows that $E_6$ has
exactly one zero in the disc. Hence $t-1$ has two zeros there,
with multiplicity, and another application of Rouch\'e shows
that $t-t_0$ has two for each $|t_0-1|\le1/500$.
This argument proves coverage of the $t$-disc, rather than
inferring coverage merely from a boundary bound. Elliptic
curves over $\C$ with the same $j$ are isomorphic. In terms of
lattices the isomorphism is a homothety; it carries the two
zeros of $\wp$ to the two zeros of the scaled $\wp$ and
preserves $N$. Statement \ref{ari:C.8.2} applies to either zero.
\end{proof}

\begin{proposition}[the first three normalized duplication steps]\label{ari:C.8.4}
Put $u=t-1$. For $n=2^k$, $k=1,2,3$, and $nP\ne O$,
\begin{equation}\label{ari:eq:ca-ladder}
 \left(\frac{N(P)}{\sqrt2}\right)^{n^2}
 =\frac{3^{3/4}}{2n}\,
       \frac{N(nP)|u|^{-1/4}}{|A_k(u)|},
\end{equation}
where
\begin{align*}
 A_1(u)&=1,\qquad A_2(u)=1-5u/27,\\
 A_3(u)&=1-\frac{217}{27}u-\frac{20486}{729}u^2
       +\frac{10910}{19683}u^3
       +\frac{161383}{177147}u^4
       -\frac{245}{14348907}u^5.
\end{align*}
On $|u|\le1/28$ we have $|A_3(u)|\ge1-10|u|$.
For every $k\ge1$, not just these first three steps, define
\[
 E_k(t)=\frac{|\Delta_E|^{(4^k-1)/12}|t-1|^{1/4}}
 {2^{4^k/2}|\psi_{2^k}(P)|}.
\]
Then $(N(P)/\sqrt2)^{4^k}=E_k(t)N(2^kP)|t-1|^{-1/4}$,
and the removable limiting value is
\[
 \lim_{t\to1}E_k(t)=\frac{3^{3/4}}{2^{k+1}}
                 =2^{1-k}\frac{3^{3/4}}4.
\]
\end{proposition}
\Uses{the hypotheses, notation and definitions stated here; \cref{ari:C.4.2}; \cref{ari:eq:cs-division}; \cref{ari:C.3.1}}

\begin{proof}
Substitute the exact division polynomials from Statement \ref{ari:C.4.2}
in \eqref{ari:eq:cs-division}. Equivalently, remove
$t^{2(n^2-1)/3}(t-1)^{n^2/4}$ from $\psi_n(P)$ and divide
the resulting polynomial by its value at $t=1$. For $n=2,4,8$
this gives the displayed polynomials and the constant
$3^{3/4}/(2n)$. All operations are rational polynomial operations;
\path{code/C_ladder_algebra.py} reproduces them and records
the next two steps as well. The general identity is \ref{ari:C.4.2}.
For its limit write $\mu=(1-\tau)/2+h$ as in \ref{ari:C.3.1}. Since
$n=2^k$ is even, $n\mu$ is congruent to $nh$ modulo the
period lattice. Hence
\[
 N(nP)=2\pi|\eta(i)|^2 n|h|(1+o(1)),\qquad
 t-1=\tfrac94\pi^4E_4(i)h^4(1+o(1)).
\]
The identity $E_4(i)^3=1728\eta(i)^{24}$ then gives
$N(nP)|t-1|^{-1/4}\to2n/3^{3/4}$.
Since the left side of the general ladder tends to one,
this proves the stated limit for every $k$.
The triangle inequality gives
$|A_3(u)|\ge1-|u|M$, where
\[
 M\le\frac{81}{10}+\frac{29}{28}
                +\frac1{28^2}+\frac1{28^3}+\frac1{28^4}<10.
\]
Here we bounded respectively the five absolute coefficients by
$81/10,29,1,1,1$. This is a rational inequality on one disc,
not a subdivision or a root search.
\end{proof}

\begin{proof}[Division-ladder proof of Theorem~\ref{fam:st:C.2.3}]
Reduce to the standard modular region. The portion $T\ge2$
is covered by Statement \ref{ari:C.8.1}. Suppose $T\le2$ and put
$r=|t-1|$. For $r\ge9/10$, the step $n=2$ in
\eqref{ari:eq:ca-ladder} and $N(2P)<17/10$ give strict inequality,
because
\[
 \frac{27(17/10)^4}{256}<\frac9{10}.
\]
There is no $2$-torsion point of the section on the smooth base.

If $r\le9/10$, then $T<7/5$. Indeed, $T\ge7/5$ would give
$|q|<1/6500$ and
\[
 |E_4|>24/25,\qquad |\Delta_0|<1/6400,
 \qquad |t|>\frac{6400(24/25)^3}{1728}>19/10,
\]
contrary to $|t|\le19/10$. The estimates follow from
\[
 \sum_{n\ge1}\sigma_3(n)r^n
 \le\frac{r(1+4r+r^2)}{(1-r)^5},\qquad
 |\Delta_0|\le r\exp\frac{24r}{1-r}.
\]
For $1/28\le r\le9/10$, the $n=4$ step and
$N(4P)<3/2$ now apply. The function
$(1-5r/27)r^{1/4}$ is increasing up to $r=27/25$.
At its left endpoint its fourth power is
\[
 \frac1{28}\left(\frac{751}{756}\right)^4
 >\frac{2187}{65536}
 =\left(\frac{3^{3/4}}8\frac32\right)^4.
\]
This proves strict inequality on the whole annulus. The lower
bound for $A_2$ is positive there, so no division by a zero has
occurred.

If $r\le1/28$, then $T<11/10$. For otherwise $|q|<1/1000$,
and the Lambert series with
\[
 \sum_{n\ge1}\sigma_5(n)r^n
 \le\frac{r(1+26r+66r^2+26r^3+r^4)}{(1-r)^7}
\]
gives $|E_6|>9/20$, $|\Delta_0|<1/900$, and hence
$|t-1|>900(9/20)^2/1728>1/28$.
On $1/500\le r\le1/28$, the $n=8$ step,
$N(8P)<143/100$, and $|A_3|\ge1-10r$ prove the claim.
Indeed $(1-10r)r^{1/4}$ has a single critical point in this
interval, a maximum at $r=1/50$, so its minimum is at an endpoint.
At both endpoints its fourth power exceeds $9/5000$, whereas
\[
 \left(\frac{3^{3/4}}{16}\frac{143}{100}\right)^4
 <\frac7{4000}<\frac9{5000}.
\]
Again the displayed positive bound excludes a division-polynomial
zero. The remaining disc $r\le1/500$ is covered by Statement
\ref{ari:C.8.3}. The smooth limiting fibre at $t=0$ has
$N=3^{1/4}<\sqrt2$; the formula also extends there by
continuity, so the singular rational model causes no omitted
modulus. The only equality occurs at $t=1$.
\end{proof}

Theorems~\ref{fam:src:C.2.3} and \cref{push:src:E.4.1,push:src:E.4.2} establish this inequality by
computer-assisted proofs. Statements \ref{ari:C.8.1}--\ref{ari:C.8.5} give another
proof with a three-step division ladder and the explicit radius
$1/50$ in the $\tau$-plane. The programs
\path{code/C_ladder_estimates.py} and
\path{code/C_local_analytic_bounds.py} record numerical diagnostics
of the constants; the proof uses the displayed rational estimates
and their series remainder bounds alone.



